\documentclass[11pt, answers]{exam}
\usepackage{tikz}
\usetikzlibrary{decorations.pathreplacing}
\usepackage{amsmath}
\usepackage{amsthm}
\usepackage[boxed]{algorithm}
\usepackage[noend]{algpseudocode}
\usepackage{babel}
\usepackage{titlesec}
\usepackage{changepage}
\usepackage{amsfonts}
\usepackage{amssymb}
\usepackage{commath}
\usepackage{mathrsfs}
\usepackage{fancybox}
\usepackage{xcolor}
\usepackage{enumitem}
\usepackage[noend]{algpseudocode}
\usepackage[margin = 1in]{geometry}
\usepackage[utf8]{inputenc}
\usepackage{pgfplots}
\pgfplotsset{compat=1.18}
\usepackage{graphicx}
\usetikzlibrary{calc}
\newif\ifsteps
\stepstrue
%\stepsfalse

\renewcommand{\solutiontitle}{\noindent}
\renewcommand{\qedsymbol}{$\blacksquare$}
\newcommand{\x}{\times}
\newcommand{\ra}{\rightarrow}
\newcommand{\imp}{\implies}
\newcommand{\ex}{\exists}
\newcommand{\prob}{\textbf{P}}
\newcommand{\R}{\mathbb{R}}
\newcommand{\C}{\mathbb{C}}
\newcommand{\F}{\mathbb{F}}
\newcommand{\N}{\mathbb{N}}
\newcommand{\Z}{\mathbb{Z}}
\newcommand{\Q}{\mathbb{Q}}
\newcommand{\comp}{\mathsf{c}}
\newcommand{\topo}{\mathcal{T}}
\newcommand{\vter}{\mathit{v_{\text{ter}}}}
\newcommand{\vex}{\mathit{v_{\text{ex}}}}
\newcommand{\lang}{\mathcal{L}}
\newcommand{\ham}{\mathcal{H}}
\newcommand{\ess}{\mathcal{S}}
\newcommand{\avg}[1]{\langle {#1} \rangle}
\newcommand{\bigavg}[1]{\left\langle {#1} \right\rangle}
\newcommand{\bracket}[2]{\langle {#1} \mid {#2}\rangle}
\newcommand{\ket}[1]{\left\lvert {#1}\right\rangle}
\newcommand{\allint}{\int_{-\infty}^{+\infty}}
\newcommand\todo[1]{\textcolor{red}{#1}}
\newcommand{\fullwidthseed}{\par\noindent\rule{\linewidth}{0pt}\vspace*{-\topskip}}
\newcommand{\vect}[1]{|{#1}\rangle}
\newcommand{\schro}{Schr{\"o}dinger}
\allowdisplaybreaks
\usepackage{calligra}
\DeclareMathAlphabet{\mathcalligra}{T1}{calligra}{m}{n}
\DeclareFontShape{T1}{calligra}{m}{n}{<->s*[2.2]callig15}{}
\newcommand{\scriptr}{\mathcalligra{r}\,}
\newcommand{\boldscriptr}{\pmb{\mathcalligra{r}}\,}
\usepackage{comment}
\includecomment{steps}

\begin{document}
	\begin{center}
		\Large \bf Useful Tips \\ Sharon Ye
	\end{center}

    \section{Trigonometric Identities}
    Instead of memorizing all the trigonometric identities, you can derive them from
    \[
    e^{i\theta}=\cos\theta+i\sin\theta \qquad\text{and}\qquad
    e^{i\alpha}e^{i\beta}=e^{i(\alpha+\beta)}
    \]
    From these two expressions we get that
    \[
    e^{ix}=\cos x+i\sin x\qquad\text{and}\qquad
    e^{-ix}=\cos x-i\sin x.
    \]
    $\implies$
    \[
    \boxed{\cos x=\frac{e^{ix}+e^{-ix}}{2}}\qquad\text{and}\qquad
    \boxed{\sin x=\frac{e^{ix}-e^{-ix}}{2i}}
    \]
    \subsection{Addition and Subtraction Formulas}
    \[
    e^{i(a+b)}=e^{ia}e^{ib}.
    \]
    Using Euler's formula,
    \begin{align*}
        \cos(a+b)+i\sin(a+b) &=(\cos a+i\sin a)(\cos b+i\sin b) \\
        &=\cos a\cos b-\sin a\sin b+i(\sin a\cos b+\cos a\sin b)
    \end{align*}
    Now equate real and imaginary parts:
    \[
    \boxed{\cos(a+b)=\cos a\cos b-\sin a\sin b}
    \]\[
    \boxed{\sin(a+b)=\sin a\cos b+\cos a\sin b}.
    \]To get subtraction formulas, simply replace \(b\to-b\):
    \[
    \boxed{\cos(a-b)=\cos a\cos b+\sin a\sin b}
    \]\[
    \boxed{\sin(a-b)=\sin a\cos b-\cos a\sin b}.
    \]

    \subsection{Double/Half Angle Forumlas}
    \[
    e^{2ix}=(e^{ix})^2.
    \]
    $\implies$
    \begin{align*}
        \cos 2x+i\sin 2x &= (\cos x+i\sin x)^2 \\
        &=\cos^2x-\sin^2x+2i\sin x\cos x
    \end{align*}
    $\implies$
    \[
    \boxed{\cos 2x=\cos^2x-\sin^2x}\qquad,\qquad
    \boxed{\sin 2x=2\sin x\cos x}
    \]
    Using the Pythatgorean identity
    \[
    1=e^{ix}e^{-ix}=(\cos x+i\sin x)(\cos x-i\sin x)
    =
    \cos^2x+\sin^2x.
    \]
    and combining it with
    \[
    \cos 2x=\cos^2x-\sin^2x
    \]
    we get the other forms of the double-angle identity:
    \[
    \boxed{\cos2x=2\cos^2x-1}\qquad,\qquad
    \boxed{\cos2x=1-2\sin^2x}
    \]Solving these gives the half-angle/power-reduction formulas:
    \[
    \boxed{\cos^2x=\frac{1+\cos2x}{2}},
    \qquad
    \boxed{\sin^2x=\frac{1-\cos2x}{2}}.
    \]
    
    \subsection{Product to Sum Identities}
    Use the exponential representation e.g.
    \begin{align*}
        \cos a\cos b &= \frac{e^{ia}+e^{-ia}}{2}\frac{e^{ib}+e^{-ib}}{2} \\
        &=\frac14
        \left(
        e^{i(a+b)}
        +e^{i(a-b)}
        +e^{-i(a-b)}
        +e^{-i(a+b)}
        \right) \\
        &=\frac14
        \left[
        2\cos(a+b)+2\cos(a-b)
        \right]
    \end{align*}
    $\implies$
    \begin{align*}
        \boxed{
        \cos a\cos b
        =
        \frac12[\cos(a+b)+\cos(a-b)]
        }
    \end{align*}
    The same procedure gives
    \[
    \boxed{
    \sin a\sin b
    =
    \frac12[\cos(a-b)-\cos(a+b)]
    }\qquad,\qquad
    \boxed{
    \sin a\cos b
    =
    \frac12[\sin(a+b)+\sin(a-b)]
    }.
    \]

    \subsection{Multiple Angle Identities and de Moivre's Formula}
    \[
    (e^{ix})^n=e^{inx} \implies
    \boxed{
    (\cos x+i\sin x)^n
    =
    \cos(nx)+i\sin(nx)
    }.
    \]This is de Moivre's formula.
    For example, for \(n=3\),
    \[
    (\cos x+i\sin x)^3
    =
    \cos3x+i\sin3x.
    \]Expand the left side. Its real part is
    \[
    \cos^3x-3\cos x\sin^2x.
    \]Thus
    \[
    \cos3x
    =
    \cos^3x-3\cos x\sin^2x.
    \]Using \(\sin^2x=1-\cos^2x\),
    \[
    {\cos3x=4\cos^3x-3\cos x}.
    \]The imaginary part similarly gives
    \[
    {\sin3x=3\sin x-4\sin^3x}.
    \]So all the \(n\)-angle identities come from one equation.
    
    \section{Binomial and Taylor Expansions}
    
    \section{Einstein Notation}
\end{document}